% ============================================================
% Capítulo 9: Documentação Viva
% ============================================================
\chapter{Documentação Viva}

\section{Introdução}

% [Placeholder: O paradoxo da documentação]
% - Ninguém gosta de escrever, todos gostam de ler
% - Documentação desatualizada é pior que nenhuma documentação
% - Documentação como produto vivo

\section{Tipos de Documentação}

% [Placeholder: Categorias essenciais]
% - Documentação de arquitetura (ADRs, diagramas)
% - README e getting started
% - Documentação de API
% - Comentários no código (quando e quando não)
% - Runbooks e playbooks de operações
% - Decision logs

\section{Documentação de Arquitetura}

% [Placeholder: Documentando decisões]
% - C4 Model (Context, Containers, Components, Code)
% - Diagramas que importam
% - ADRs como documentação viva
% - Ferramentas: Structurizr, PlantUML, Mermaid

\section{Documentação como Código}

% [Placeholder: Docs as code]
% - Markdown no repositório
% - Versionamento junto com o código
% - Review de documentação no PR
% - Geração automática de docs

\section{IA e Documentação}

% [Placeholder: Automatizando com IA]
% - Geração automática de README
% - Explicação de código complexo
% - Geração de documentação de API
% - Atualização de docs após mudanças
% - Os limites da IA na documentação

\section{Documentação de APIs}

% [Placeholder: APIs bem documentadas]
% - OpenAPI/Swagger
% - Documentação de GraphQL
% - Exemplos de uso
% - Versionamento de APIs
% - Ferramentas: Redocly, Postman

\section{Manutenção da Documentação}

% [Placeholder: Mantendo docs vivas]
% - Checklist de atualização em PRs
% - Links quebrados e docs órfãs
% - Feedback loop com leitores
% - Métricas de uso de documentação
% - Cultura de documentação na equipe

\section{O Mínimo Viável de Documentação}

% [Placeholder: Quanto文档 é suficiente]
% - README obrigatório
% - Arquitetura em 1 página
% - Como rodar o projeto localmente
% - Como fazer deploy
% - Onde encontrar ajuda

\section{Conclusão}

% [Placeholder: Resumo]
% - Documentação é ato de empatia
% - Pouca, boa e atualizada > muita, ruim e esquecida

\section*{Exercícios}

% [Placeholder: Exercícios]
% 1. Escreva um README para um projeto existente usando IA.
%    Compare com um README escrito manualmente.
% 2. Crie um ADR para uma decisão recente em seu projeto.
% 3. Audit a documentação de um projeto. Identifique o que está
%    desatualizado.
